\subsection{Kirby calculus for closed oriented marked four-manifolds}
\label{subsec:marked-kirby-calculus}

The unmarked completeness theorem,
\cref{prop:special-handle-movie-completeness}, already supplies the required
Kirby movie.  We only have to lift that movie through a prescribed class in
second homology.  The lift is forced: represent the Poincar\'e dual of the
class in handle cohomology, and transport it by the diffeomorphism carried
by each move.  We make this precise and record the resulting coordinate
rules needed later.

\begin{definition}[Marked four-manifold]
\label{def:marked-four-manifold}
A \emph{marked four-manifold} is a pair \((X,\alpha)\), where \(X\) is a
closed, connected, oriented smooth four-manifold and
\(\alpha\in H_2(X;\mathbb Z)\).
\end{definition}

\begin{definition}[Handle cohomology]
\label{def:handle-cohomology}
For a special Kirby presentation \(\mathcal H\), let
\[
  C^k(\mathcal H):=\operatorname{Hom}(C_k(\mathcal H),\mathbb Z),
  \qquad
  \delta^k(\lambda):=\lambda\circ\partial_{k+1}.
\]
We write
\[
  Z^2(\mathcal H):=\ker(\delta^2),
  \qquad
  H^2(\mathcal H):=
  \ker(\delta^2)/\operatorname{im}(\delta^1).
  \tag{3.2}\label{eq:handle-cohomology}
\]
\end{definition}

\begin{proposition}[Handle-cohomology realization]
\label{prop:handle-homology-realization}
If \(\mathcal H\) presents \(X\), cellular-to-singular comparison gives a
natural isomorphism
\[
  \kappa_{\mathcal H}\colon
  H^2(\mathcal H)\xrightarrow{\ \cong\ }H^2(X;\mathbb Z).
  \tag{3.3}\label{eq:handle-comparison}
\]
\end{proposition}

\begin{proof}
The handle filtration is a CW structure on \(X\).  Its cellular boundary is
the signed attaching-sphere--belt-sphere intersection map, hence its
cellular chain complex is the handle complex
\eqref{eq:handle-complex}.  The claim is therefore the natural
cellular-to-singular comparison isomorphism; see
\cite[Section~4.2]{GompfStipsicz1999} and
\cite[Section~3]{DurstGeigesKegel2019}.
\end{proof}

\begin{definition}[Marked special Kirby presentation]
\label{def:marked-special-kirby-presentation}
\label{def:handle-homology-marking}
Let \((X,\alpha)\) be a marked four-manifold and let \(\mathcal H\) be a
special Kirby presentation of \(X\).  Its marking is the uniquely determined
class
\[
  \xi_{\mathcal H}
  :=\kappa_{\mathcal H}^{-1}\bigl(\operatorname{PD}_X(\alpha)\bigr)
  \in H^2(\mathcal H).
  \tag{3.4}\label{eq:marked-presentation}
\]
The pair \((\mathcal H,\xi_{\mathcal H})\) is a
\emph{marked special Kirby presentation} of \((X,\alpha)\).

A \emph{handle cocycle} is a representative
\(a\in Z^2(\mathcal H)\) of \(\xi_{\mathcal H}\).  If \(e_j\) is the
oriented basis element of \(C_2(\mathcal H)\) represented by the \(j\)th
framed component, write
\[
  a_j:=a(e_j)\in\mathbb Z.
  \tag{3.5}\label{eq:handle-marking-coordinates}
\]
The cocycle \(a\) is auxiliary data; the marking is its cohomology class.
\end{definition}

\begin{definition}[Canonical marking transport]
\label{def:canonical-marking-transport}
Let \(m\colon\mathcal H\to\mathcal H'\) be one of the elementary moves in
\cref{def:kirby-move}, or a \(3/3\) slide.  Its standard local model
determines a relative isotopy class of orientation-preserving
diffeomorphisms
\[
  [\phi_m]\colon X_{\mathcal H}\longrightarrow X_{\mathcal H'},
\]
equal to the identity outside the support of the move.  Define
\[
  T_m:=
  \kappa_{\mathcal H'}^{-1}
  \circ(\phi_m^{-1})^*
  \circ\kappa_{\mathcal H}
  \colon H^2(\mathcal H)\longrightarrow H^2(\mathcal H').
  \tag{3.6}\label{eq:marking-transport}
\]
This is independent of the representative of \([\phi_m]\).
\end{definition}

\begin{proposition}[Handle-chain description of the transport]
\label{lem:handle-cohomology-naturality}
Every move \(m\) above induces a chain-homotopy equivalence
\[
  F_m\colon C_*(\mathcal H)\longrightarrow C_*(\mathcal H').
\]
For the dual cochain map
\[
  F_m^\vee\colon C^*(\mathcal H')\longrightarrow C^*(\mathcal H),
  \qquad b\longmapsto b\circ F_m,
\]
the comparison square
\[
  \kappa_{\mathcal H}\circ H^2(F_m^\vee)
  =\phi_m^*\circ\kappa_{\mathcal H'}
  \tag{3.7}\label{eq:single-move-comparison-square}
\]
commutes.  Consequently,
\[
  T_m=H^2(F_m^\vee)^{-1}.
  \tag{3.8}\label{eq:canonical-transport-chain-map}
\]
For a standard \(k\)--\((k+1)\) birth, \(k=1,2\), the target complex is
\[
  C_*(\mathcal H')\cong C_*(\mathcal H)\oplus A_*,
  \qquad
  A_{k+1}=\mathbb Zu,\quad A_k=\mathbb Zv,\quad
  \partial u=\varepsilon v,\quad\varepsilon\in\{\pm1\}.
  \tag{3.9}\label{eq:acyclic-birth-summand}
\]
\end{proposition}

\begin{proof}
The handle groups are the relative cellular groups of the handle
filtration.  Naturality of cellular-to-singular comparison therefore proves
\eqref{eq:single-move-comparison-square} once the cellular map is identified.
There are only three cases.

\begin{enumerate}[label=(\roman*),ref=\roman*]
  \item A level isotopy, including a \(2/1\) move, induces the cellular
  continuation map.

  \item For a same-index \(k/k\) slide, let
  \(J_m\colon C_*(\mathcal H')\to C_*(\mathcal H)\) send the new moving
  generator to the old moving generator plus or minus the old stationary
  generator and fix the others.  The post-slide boundary matrix is the
  pre-slide matrix in this new basis, so \(J_m\partial'=\partial J_m\);
  take \(F_m=J_m^{-1}\).

  \item A birth is the elementary expansion
  \eqref{eq:acyclic-birth-summand}, with \(F_m\) the inclusion of the old
  summand; a death is the inverse collapse.  The added summand is
  contractible.
\end{enumerate}

These are precisely the cellular maps of the standard local models
\cite[Sections~4.1--4.2]{DurstGeigesKegel2019}.  Thus the comparison square
commutes in every case.  Since \(F_m\) is a chain-homotopy equivalence,
inverting that square gives \eqref{eq:canonical-transport-chain-map}.
\end{proof}

\begin{definition}[Marked Kirby move]
\label{def:marked-kirby-move}
A \emph{marked Kirby move}
\[
  m\colon(\mathcal H,\xi)\longrightarrow(\mathcal H',\xi')
\]
is an elementary Kirby move, or a \(3/3\) slide, satisfying
\[
  \xi'=T_m\xi.
  \tag{3.10}\label{eq:marked-move-condition}
\]
\end{definition}

\begin{proposition}[Coordinate form of marking transport]
\label{prop:marking-coordinate-transport}
Let \(a\in Z^2(\mathcal H)\) represent \(\xi\).  The transported class
\(T_m\xi\) has a cocycle representative \(a'\) with the following
coordinates.
\begin{enumerate}[label=(\roman*),ref=\roman*]
  \item For framed isotopy, a \(1/1\) slide, a \(2/1\) slide, or a
  \(3/3\) slide, the \(2\)-handle generators are canonically identified and
  \(a'=a\).

  \item For a \(2/2\) slide, let \(e_x\) be the moving generator,
  \(e_y\) the stationary generator, and choose
  \(\varepsilon\in\{\pm1\}\) so that
  \(J_m(e_x')=e_x+\varepsilon e_y\).  Then
  \[
    a_x'=a_x+\varepsilon a_y,
    \qquad a_y'=a_y,
    \tag{3.11}\label{eq:two-two-marking-transport}
  \]
  and all uninvolved coordinates are unchanged.

  \item A \(1\)--\(2\) birth admits the representative
  \(a'=a\oplus0\).  Every representative of the transported class is
  cohomologous to one whose new \(2\)-handle coordinate is zero.

  \item For a \(2\)--\(3\) birth, the cocycle equation forces the new
  \(2\)-handle coordinate to be zero; hence again \(a'=a\oplus0\).
\end{enumerate}
The death formulas are the inverses of the corresponding birth formulas.
\end{proposition}

\begin{proof}
By \cref{lem:handle-cohomology-naturality}, these statements are obtained by
dualizing the three chain maps in its proof.  Only the last three
calculations require comment.  For a \(2/2\) slide, take \(a'=a\circ J_m\);
then \(F_m^\vee(a')=a\), and evaluation on \(e_x'\) gives
\[
  a_x'=a(e_x+\varepsilon e_y)=a_x+\varepsilon a_y.
\]
For a \(1\)--\(2\) birth, the new degree-two cochain lies in
\(\operatorname{im}(\delta^1)\), so its coordinate can be removed by a
coboundary.  For a \(2\)--\(3\) birth, the new component of
\(\delta^2\) is multiplication by \(\varepsilon\), so the cocycle equation
forces that coordinate to vanish.
\end{proof}

\begin{lemma}[Naturality along a handle movie]
\label{lem:marking-naturality-along-movie}
For a finite movie
\[
  M=(m_1,\ldots,m_N)\colon\mathcal H_0\longrightarrow\mathcal H_N,
  \qquad
  T_M:=T_{m_N}\circ\cdots\circ T_{m_1},
\]
with realized diffeomorphism \(\Phi_M\colon X_0\to X_N\), one has
\[
  \kappa_{\mathcal H_N}\circ T_M
  =(\Phi_M^{-1})^*\circ\kappa_{\mathcal H_0}.
  \tag{3.12}\label{eq:movie-marking-naturality}
\]
\end{lemma}

\begin{proof}
Compose \eqref{eq:marking-transport} along the finite movie.  The
contravariance of cohomology reverses the order exactly as displayed.
\end{proof}

\begin{theorem}[Marked Kirby--Cerf completeness]
\label{prop:marked-kirby-cerf}
Let \(f\colon X\to Y\) be an orientation-preserving diffeomorphism of closed,
connected, oriented smooth four-manifolds, and suppose
\(f_*\alpha=\alpha'\).  If
\((\mathcal H,\xi_{\mathcal H})\) and
\((\mathcal H',\xi_{\mathcal H'})\) present
\((X,\alpha)\) and \((Y,\alpha')\), respectively, then the unmarked movie of
\cref{prop:special-handle-movie-completeness} lifts uniquely to a finite
marked movie
\[
  (\mathcal H,\xi_{\mathcal H})
  \rightsquigarrow
  (f^*\mathcal H',\widetilde\xi_{\mathcal H'})
\]
on \(X\).  Here
\(\widetilde\xi_{\mathcal H'}\) is the marking of the pulled-back
presentation.  Following this movie by its tautological endpoint
identification with \((\mathcal H',\xi_{\mathcal H'})\) represents the
prescribed isotopy class of \(f\).
\end{theorem}

\begin{proof}
Take the unmarked movie
\[
  \mathcal H=\mathcal H_0
  \xrightarrow{m_1}\cdots\xrightarrow{m_N}
  \mathcal H_N=f^*\mathcal H'
\]
given by \cref{prop:special-handle-movie-completeness}.  Starting from
\(\xi_0=\xi_{\mathcal H}\), define recursively
\(\xi_r:=T_{m_r}\xi_{r-1}\).  By
\eqref{eq:marked-move-condition}, every step is then a marked move, and this
is the only possible lift.

The realized diffeomorphism of the fixed-manifold movie is isotopic to
\(\operatorname{id}_X\).  Hence \eqref{eq:movie-marking-naturality} gives
\[
  \kappa_{\mathcal H_N}(\xi_N)
  =\operatorname{PD}_X(\alpha).
\]
By the uniqueness in \eqref{eq:marked-presentation},
\(\xi_N=\widetilde\xi_{\mathcal H'}\).

Finally, naturality of Poincar\'e duality gives
\[
  f^*\operatorname{PD}_Y(\alpha')
  =f^*\operatorname{PD}_Y(f_*\alpha)
  =\operatorname{PD}_X(\alpha).
  \tag{3.13}\label{eq:pd-marking-naturality}
\]
Thus the tautological endpoint identification induced by \(f\) carries the
pulled-back marking to \(\xi_{\mathcal H'}\).  The corresponding unmarked
composite represents \(f\) by
\cref{prop:special-handle-movie-completeness}, completing the proof.
\end{proof}
